\section{Strategy}
\label{sec:pat-strategy}

\problemstatement{Define a family of interchangeable algorithms,
encapsulate each one, and make them swappable at run time so the algorithm
varies independently of the client that uses it.}

\begin{patternbox}[When You Reach for It]
The trigger is \emph{``several ways to do the same thing, chosen at run
time''} -- different sort orders, payment methods, compression schemes,
route-finding algorithms. When you catch yourself writing \texttt{if
(mode == A) doA(); else doB();}, extract the modes into strategy objects.
\end{patternbox}

\subsection*{Understanding the pattern}

Encapsulate each algorithm behind a common interface. The context holds a
pointer to a strategy and delegates the varying step to it, without knowing
which concrete algorithm it is. Clients pick and inject a strategy; they
can change it at run time. This is the Open/Closed Principle for
behaviour: add a new algorithm as a new class, edit nothing existing.

\begin{ideabox}[Inject the Algorithm as an Object]
Pull the part that varies into a strategy object and hand it to the
context. The context is written once against the strategy interface; the
concrete behaviour is a plug-in chosen by the client. This is composition
replacing a conditional.
\end{ideabox}

\subsection*{C++ implementation}

\begin{lstlisting}
#include <bits/stdc++.h>
using namespace std;

struct PaymentStrategy {
    virtual void pay(int amount) = 0;
    virtual ~PaymentStrategy() = default;
};
struct CreditCard : PaymentStrategy {
    void pay(int amt) override { cout << "Paid " << amt << " by card\n"; }
};
struct UpiPayment : PaymentStrategy {
    void pay(int amt) override { cout << "Paid " << amt << " via UPI\n"; }
};

class Checkout {                                // context
    unique_ptr<PaymentStrategy> strategy;
public:
    void setStrategy(unique_ptr<PaymentStrategy> s) { strategy = move(s); }
    void pay(int amount) { strategy->pay(amount); }   // delegate varying step
};

int main() {
    Checkout checkout;
    checkout.setStrategy(make_unique<CreditCard>());
    checkout.pay(500);
    checkout.setStrategy(make_unique<UpiPayment>());  // swap at run time
    checkout.pay(300);
    return 0;
}
\end{lstlisting}

\begin{complexitybox}[Trade-offs]
\textbf{Pros.} Algorithms are interchangeable at run time; removes
conditionals; each algorithm is isolated and testable; supports
Open/Closed. \textbf{Cons.} Clients must know the strategies exist to
choose one; more objects/classes; overkill when the behaviour never varies.
\end{complexitybox}

\begin{takeawaybox}
Strategy encapsulates interchangeable algorithms behind one interface and
lets the client inject the choice, replacing behavioural conditionals with
composition. It is the single most common behavioural pattern in
interviews. Contrast with State: same shape, but the \emph{client} picks a
Strategy while a State machine swaps its own state.
\end{takeawaybox}
